\documentclass[10pt]{article}

% ---------- Document Identity (update these only) ----------
\newcommand{\DocStage}{}
\newcommand{\DocTitle}{Your Road to Tier One \textbf{SOF}}
\newcommand{\ProgramName}{182-Week Civilian-to-Selection Readiness Pipeline}
\newcommand{\ProgramSubtitle}{}

% ---------- Page ----------
\usepackage[legalpaper,left=0.5in,right=0.5in,top=0.6in,bottom=0.45in]{geometry}

% ---------- Typography ----------
\usepackage[T1]{fontenc}
\usepackage[utf8]{inputenc}
\usepackage{libertinus}
\usepackage{microtype}
\usepackage{parskip}
\usepackage{ragged2e}
\usepackage{textcomp}
\AtBeginDocument{\color{black}}
\AtBeginDocument{\pagecolor{white}}
\AtBeginDocument{\justifying}

% ---------- Landscape pages ----------
\usepackage{pdflscape}

% ---------- Color ----------
\usepackage[table]{xcolor}
\definecolor{StageBlue}{HTML}{1F4E79}
\definecolor{StageBlueDark}{HTML}{163A5A}
\definecolor{LightGray}{HTML}{F2F4F7}
\definecolor{MidGray}{HTML}{D9DEE5}
\definecolor{RuleGray}{HTML}{C7CED8}
\definecolor{HeaderInk}{HTML}{000000}
\definecolor{Ink}{HTML}{000000}

% ---------- Utility ----------
\usepackage{enumitem}
\sloppy
\setlength{\emergencystretch}{2em}

% ---------- Boxes / UI ----------
\usepackage[most]{tcolorbox}

\tcbset{
  charterTitle/.style={
    enhanced,
    colback=StageBlue!4,
    colframe=StageBlue,
    boxrule=1.25pt,
    arc=3pt,
    left=16pt,right=16pt,top=14pt,bottom=14pt,
    borderline north={3.2pt}{0pt}{StageBlue},
    borderline west={4pt}{0pt}{StageBlue},
    borderline south={1.25pt}{0pt}{StageBlue},
    drop shadow={black!25!white},
  },
  charterRule/.style={
    enhanced,
    colback=LightGray,
    colframe=RuleGray,
    boxrule=0.85pt,
    arc=3pt,
    left=12pt,right=12pt,top=10pt,bottom=10pt,
    borderline west={3pt}{0pt}{StageBlue},
  },
  charterBadge/.style={
    enhanced,
    colback=StageBlue,
    coltext=white,
    colframe=StageBlueDark,
    boxrule=0.6pt,
    arc=2pt,
    left=6pt,right=6pt,top=3pt,bottom=3pt,
  }
}
\newtcbox{\Badge}{nobeforeafter,charterBadge}

% ---------- Lists ----------
\setlist[itemize]{leftmargin=1.5em, itemsep=0.25em, topsep=0.25em}
\setlist[enumerate]{leftmargin=1.8em, itemsep=0.25em, topsep=0.25em}

% ---------- TOC ----------
\usepackage{tocloft}
\setcounter{tocdepth}{3}
\setlength{\cftbeforesecskip}{3pt}
\setlength{\cftbeforesubsecskip}{2pt}
\setlength{\cftbeforesubsubsecskip}{0pt}
\renewcommand{\cftsecfont}{\bfseries}
\renewcommand{\cftsecpagefont}{\bfseries}
\renewcommand{\cftsubsecfont}{\normalfont}
\renewcommand{\cftsubsecpagefont}{\normalfont}
\renewcommand{\cftsubsubsecfont}{\normalfont}
\renewcommand{\cftsubsubsecpagefont}{\normalfont}
\setlength{\cftsecindent}{0pt}
\setlength{\cftsubsecindent}{1.25em}
\setlength{\cftsubsubsecindent}{2.8em}
\setlength{\cftsecnumwidth}{2.1em}
\setlength{\cftsubsecnumwidth}{2.8em}
\setlength{\cftsubsubsecnumwidth}{3.6em}

% Remove the built-in TOC heading so only the custom heading is shown
\makeatletter
\renewcommand{\tableofcontents}{\@starttoc{toc}}
\makeatother

% ---------- Header/Footer: page number only ----------
\usepackage{fancyhdr}
\pagestyle{fancy}
\fancyhf{}
\renewcommand{\headrulewidth}{0pt}
\renewcommand{\footrulewidth}{0pt}
\fancyfoot[C]{\thepage}

% ---------- Section formatting ----------
\usepackage{titlesec}

\titleformat{\section}
  {\Large\bfseries\color{Ink}}
  {\thesection}{0.6em}{}
  [\vspace{0.25em}\textcolor{StageBlue}{\rule{\linewidth}{0.8pt}}]

\titleformat{\subsection}
  {\large\bfseries\color{Ink}}
  {\thesubsection}{0.6em}{}

\titleformat{\subsubsection}
  {\normalsize\bfseries\color{Ink}}
  {\thesubsubsection}{0.6em}{}

\titlespacing*{\section}{0pt}{0.95em}{0.35em}
\titlespacing*{\subsection}{0pt}{0.65em}{0.25em}
\titlespacing*{\subsubsection}{0pt}{0.55em}{0.2em}

% ---------- Table engine ----------
\usepackage{tabularray}
\UseTblrLibrary{booktabs}

% ---------- tabularray: GLOBAL table treatment ----------
\SetTblrInner{
  cells = {fg=black, bg=white, valign=t},
}

% ---------- Header helpers ----------
\newcommand{\Hdr}[1]{\shortstack[c]{\strut #1\strut}}

% ---------- Lines ----------
\newcommand{\TitleLine}{\textcolor{StageBlue}{\rule{0.98\linewidth}{0.95pt}}}
\newcommand{\SubTitleLine}{\textcolor{RuleGray}{\rule{0.98\linewidth}{0.65pt}}}

% ---------- Continued on next page ----------
\newcommand{\ContinuedOnNextPageInline}{%
  \par
  \vfill
  \noindent\textcolor{RuleGray}{\rule{\linewidth}{1pt}}\vspace{-2pt}
  \noindent\hfill\textit{Continued on next page}\par\vspace{-10pt}
  \noindent\textcolor{RuleGray}{\rule{\linewidth}{1pt}}
  \clearpage
}

% ---------- PDF hyperlinks ----------
\usepackage[hidelinks]{hyperref}
\hypersetup{
  colorlinks=false,
  pdfborder={0 0 0},
  linktoc=all,
  pdftitle={\DocTitle},
  pdfauthor={\ProgramName}
}

% =========================================================
\begin{document}

\section*{9.E.1 Program Header}
\phantomsection
\addcontentsline{toc}{section}{9.E.1 Program Header}

Program Name: \textbf{SOF} Physical Development Transformation Program

Version: v1.0

Date (YYYY-MM-DD): \rule{1.75in}{0.35pt}

Architect: \rule{2.25in}{0.35pt}

Based On: Stages 1–9

Notes:

\begin{itemize}
\item Release posture: Program v1.0 is a packaged, deployment-ready doctrine set. The PDF binder is canonical read-only doctrine; operational artifacts are maintained in the companion folder bundle under controlled change discipline.
\item Canonical storage pointer for versioning artifacts:
\begin{itemize}
\item /99\_Change\_Log\_and\_Archive/Versioning/
\end{itemize}
\item Stage 8 Issue IDs (\textbf{ISS-08}-\#\#\#) are authoritative reason references for known problems carried forward as Known Issues until resolved via controlled patching and revalidation.
\end{itemize}

\section*{9.E.2 Change-Log Template}
\phantomsection
\addcontentsline{toc}{section}{9.E.2 Change-Log Template}

\begingroup
\scriptsize
\setlength{\tabcolsep}{2pt}
\renewcommand{\arraystretch}{1.12}

\begin{longtblr}{
  width=\linewidth,
  rowhead=1,
  colspec = {
    Q[l,wd=0.30in]
    Q[l,wd=0.325in]
    Q[l,wd=.6in]
    Q[l,wd=.4in]
    X[l,1.55]
    X[l,1.85]
    X[l,2.10]
    Q[l,wd=0.5in]
    X[l,1.75]
    X[l,1.10]
  },
  hlines,
  vlines,
  cells = {hsep=6pt, vsep=6pt, halign=l, valign=t},
  row{odd} = {bg=white},
  row{even} = {bg=LightGray},
  row{1} = {bg=StageBlue!15, fg=black, font=\bfseries, halign=l, valign=m},
}
\Hdr{Version} &
\Hdr{Patch \textbf{ID}} &
\Hdr{Date \\YYYY-MM-DD} &
\Hdr{Changed \\By Role} &
\Hdr{Summary of \\Changes} &
\Hdr{Affected \\Stages/Files} &
\Hdr{Reason Issue \textbf{ID} \\or Observation} &
\Hdr{Impact Level \\\textbf{L1}/\textbf{L2}/\textbf{L3}} &
\Hdr{Revalidation \\Scope} &
\Hdr{Notes} \\
 & & & & & & & & & \\
\end{longtblr}

\endgroup

\section*{9.E.3 Guidance}
\phantomsection
\addcontentsline{toc}{section}{9.E.3 Guidance}

\begin{itemize}
\item Version meanings and usage conditions:
  \begin{itemize}
  \item Major (v2.0, v3.0, …): Used only when program identity, phase structure, or core governance constraints change in a way that materially alters downstream artifacts and deployment posture. Major versions require full revalidation and repackaging.
  \item Minor (v1.1, v1.2, …): Used when doctrine meaning changes across one or more stages in a material but bounded way (e.g., adjustments to phase gating language, category timing posture, or test cadence posture) that affects execution and requires crosswalk revalidation.
  \item Patch (v1.0.1, v1.0.2, …): Used for discrete corrections and controlled updates that do not change program identity and are scoped to specific artifacts/sections. Patch releases still require explicit Patch \textbf{ID}, Impact Level, and Revalidation Scope when doctrine meaning is affected.
  \end{itemize}

\item Full-phase freeze window rule:
  \begin{itemize}
  \item Freeze window runs from phase start through end-gate decision.
  \item Only safety substitutions are allowed without a patch.
  \item Any structural change (tests referenced, gate composition, cadence pattern, required capability timing, required logging fields posture) requires a patch entry, briefing, and revalidation posture per Stage 1.F Impact Level discipline.
  \end{itemize}

\item Archive rule:
  \begin{itemize}
  \item Older releases are stored only in /99\_Change\_Log\_and\_Archive/ (do not leave old versions in main stage folders).
  \item Canonical versioning artifacts live in /99\_Change\_Log\_and\_Archive/Versioning/:
    \begin{itemize}
    \item \textbf{CHANGE-LOG}
    \item \textbf{PATCH-REGISTER}
    \item \textbf{FREEZE-WINDOWS}
    \end{itemize}
  \item The PDF binder must include a pointer/index to /99\_Change\_Log\_and\_Archive/Versioning/ so implementers can locate the authoritative change-control record.
  \end{itemize}

\item Revalidation rule:
  \begin{itemize}
  \item Any upstream doctrine change that affects tests, metrics, gates, protected windows, tier timing, or category dependencies requires Stage 8 impact assessment and crosswalk re-run as applicable.
  \item Impact Level discipline:
    \begin{itemize}
    \item \textbf{L1}: wording clarity/local formatting only; revalidation scope is local artifact review.
    \item \textbf{L2}: doctrine change within a stage plus direct downstream consumers; revalidation scope includes direct downstream stages and affected crosswalks.
    \item \textbf{L3}: constraint change; revalidation scope requires full Stage 8 crosswalk re-run and updated Known Issues status before redeployment.
    \end{itemize}
  \end{itemize}

\item Single source of truth posture:
  \begin{itemize}
  \item The binder is canonical read-only doctrine used for gate decisions and audit traceability.
  \item The versioning files (\textbf{CHANGE-LOG}, \textbf{PATCH-REGISTER}, \textbf{FREEZE-WINDOWS}) are canonical in /99\_Change\_Log\_and\_Archive/Versioning/ and the binder contains a pointer to that location.
  \item No silent edits: any non-cosmetic change is recorded with Version, Patch \textbf{ID} (when applicable), Date, Changed By Role, Impact Level, and Revalidation Scope, with \textbf{ISS-08}-\#\#\# referenced when relevant.
  \end{itemize}
\end{itemize}

\end{document}